\documentclass[10pt,a4paper]{article}

\usepackage[utf8]{inputenc}
\usepackage[T1]{fontenc}
\usepackage[margin=2cm]{geometry}
\usepackage{amsmath,amssymb}
\usepackage{xcolor}
\usepackage{tcolorbox}
\tcbuselibrary{breakable}
\usepackage{hyperref}
\usepackage{enumitem}

\definecolor{partcolor}{RGB}{0,51,102}
\definecolor{chaptercolor}{RGB}{0,102,204}

\newtcolorbox{partbox}{
    colback=partcolor!5,colframe=partcolor,
    fonttitle=\bfseries\large,breakable,
    title=PARTIE
}

\newtcolorbox{chapterbox}[1]{
    colback=chaptercolor!5,colframe=chaptercolor,
    fonttitle=\bfseries,breakable,
    title=#1
}

\title{\textbf{Plan Complet du Cours}\\
\LARGE Programmation C Avancée\\
\large L3 - Université Cheikh Anta Diop de Dakar}
\author{Dr. El Hadji Bassirou TOURÉ\\
Département de Mathématiques et Informatique\\
Faculté des Sciences et Techniques}
\date{Version 1.0 - Novembre 2025}

\begin{document}

\maketitle

\tableofcontents
\newpage

\section*{Vue d'ensemble du cours}

Ce cours couvre la programmation C avancée avec une approche orientée systèmes. L'objectif est de former des programmeurs capables de comprendre les mécanismes sous-jacents et d'écrire du code robuste, performant et sécurisé.

\subsection*{Structure pédagogique}

\begin{itemize}[leftmargin=2cm]
\item \textbf{CM (Cours Magistral):} 8-12 pages, contenu technique dense avec visualisations mémoire TikZ
\item \textbf{TD (Travaux Dirigés):} 3-5 pages, exercices forensics avec dumps hexadécimaux, prédiction de sorties
\item \textbf{TD\_Corrigé:} Solutions détaillées avec explications pas-à-pas
\item \textbf{TP (Travaux Pratiques):} 4-6 pages, labs pratiques (compilation, debugging, algorithmes)
\item \textbf{TP\_Pistes:} Indices uniquement, pas de solutions complètes pour les défis algorithmiques
\end{itemize}

\subsection*{Environnement de travail}

Tous les TPs utilisent un environnement Docker standardisé:
\begin{verbatim}
Image: elbachirtoure/c-advanced-course:v1.0
Outils: gcc, gdb, valgrind, vim, make
\end{verbatim}

\subsection*{Référence principale}

\textit{Computer Systems: A Programmer's Perspective (3rd Edition)} par Randal E. Bryant \& David R. O'Hallaron

\newpage

\section{TP-1: Setup Environnement}

\textbf{Fichier unique:} \texttt{TP-1\_Setup\_Environnement.tex}

\subsection*{Objectifs}
\begin{itemize}
\item Installer et configurer Docker
\item Maîtriser les commandes Docker de base
\item Introduction à vim (modes, commandes essentielles)
\item Premier programme C et compilation
\end{itemize}

\subsection*{Contenu du TP}
\begin{itemize}
\item Installation Docker (Windows/macOS/Linux)
\item Lancement du container et volumes partagés
\item Vim: modes (normal, insertion, commande)
\item Vim: navigation, édition, sauvegarde
\item Exercices vim progressifs
\item Premier programme "Hello World"
\item Compilation avec gcc
\item Exécution du programme
\end{itemize}

\textbf{Durée estimée:} 2-3 heures

\textbf{Prérequis:} Aucun

\textbf{Livrables:} Environnement fonctionnel + exercices vim complétés

\newpage

\section{Ch-1: Introduction au Cours}

\textbf{Fichier unique:} \texttt{CM\_Ch-1\_Introduction.tex} (2-3 pages maximum)

\subsection*{Objectifs}
\begin{itemize}
\item Comprendre la pertinence du langage C aujourd'hui
\item Connaître l'histoire et l'évolution du C
\item Avoir une vue d'ensemble de la structure du cours
\end{itemize}

\subsection*{Contenu}

\textbf{1. Pourquoi le C?}
\begin{itemize}
\item Exemples concrets: Linux kernel, PostgreSQL, nginx, FFmpeg
\item TensorFlow/PyTorch backends, OpenSSL, Git
\item Systèmes embarqués, drivers, OS
\end{itemize}

\textbf{2. Histoire du C}
\begin{itemize}
\item Dennis Ritchie, Bell Labs (1972)
\item K\&R C, ANSI C (C89), C99, C11, C17, C23
\item Influence sur C++, Java, C\#, etc.
\end{itemize}

\textbf{3. Structure du cours}
\begin{itemize}
\item Les 6 grandes parties (fondamentaux, représentation, mémoire, structures, exécution, outils)
\item Objectifs généraux du cours
\item Approche pédagogique (comprendre le "comment ça marche")
\end{itemize}

\textbf{4. Prérequis}
\begin{itemize}
\item Connaissance basique du C (variables, fonctions, boucles)
\item Notions d'algorithmique
\item Motivation pour la programmation système
\end{itemize}

\textbf{5. Environnement de travail}
\begin{itemize}
\item Docker container standardisé
\item Outils: gcc, gdb, valgrind, vim
\end{itemize}

\textbf{6. Organisation pédagogique}
\begin{itemize}
\item CM/TD/TP avec exercices forensics
\item Projets pratiques
\end{itemize}

\textbf{Note:} PAS de section sur l'évaluation, l'intégrité académique, ou les conseils méthodologiques.

\newpage

\begin{partbox}
PARTIE I: FONDAMENTAUX C
\end{partbox}

\section{Ch0: Syntaxe C Complète}

\textbf{Référence:} Révision complète des bases du C

\subsection*{CM: Cours Magistral (8-10 pages)}

\textbf{Contenu théorique:}
\begin{enumerate}
\item \textbf{Structures de Contrôle}
   \begin{itemize}
   \item Blocs et portée (scope)
   \item if/else, dangling else, opérateur ternaire
   \item switch et fall-through
   \item Boucles: while, do-while, for
   \item break, continue, goto (usage légitime)
   \end{itemize}

\item \textbf{Types, Opérateurs et Conversions}
   \begin{itemize}
   \item Types primitifs et tailles (architecture-dependent)
   \item Types portables (\texttt{<stdint.h>})
   \item Opérateurs: arithmétiques, relationnels, logiques, bit-à-bit
   \item Priorité et associativité des opérateurs
   \item Conversions implicites et explicites (cast)
   \item Pièges: signed/unsigned, promotion entière
   \end{itemize}

\item \textbf{Déclaration vs Définition}
   \begin{itemize}
   \item Différence déclaration/définition
   \item One Definition Rule (ODR)
   \item Mots-clés: extern, static
   \item Erreurs de compilation vs linking
   \end{itemize}

\item \textbf{Points de Séquence et Comportements Indéfinis}
   \begin{itemize}
   \item Définition des points de séquence
   \item Liste exhaustive des points de séquence
   \item Undefined Behavior (UB): définition et dangers
   \item Exemples: modifications multiples, ordre d'évaluation
   \item Débordement d'entiers signés vs non-signés
   \item Court-circuit des opérateurs logiques
   \end{itemize}

\item \textbf{Fonctions et Variables}
   \begin{itemize}
   \item Déclaration et définition de fonctions
   \item Passage par valeur vs simulation passage par référence
   \item Paramètres const
   \item Portée: bloc, fonction, fichier, globale
   \item Classes de stockage: auto, static (local/global), const, volatile
   \item Récursivité
   \end{itemize}

\item \textbf{Compilation et Debugging}
   \begin{itemize}
   \item Les 4 phases: preprocessing, compilation, assemblage, linking
   \item Options gcc: -Wall, -Wextra, -Werror, -g, -O0/-O2/-O3
   \item GDB: commandes essentielles (break, run, next, step, print, backtrace)
   \item Introduction à Valgrind (aperçu uniquement)
   \end{itemize}
\end{enumerate}

\textbf{Visualisations:}
\begin{itemize}
\item Diagrammes de portée (scope)
\item Priorité des opérateurs (tableau)
\item Phases de compilation (schéma)
\item PAS de visualisations mémoire détaillées (Ch2+)
\end{itemize}

\subsection*{TD: Travaux Dirigés (3-4 pages)}

\textbf{Focus: Prédire la sortie et identifier les erreurs}

\textbf{Exercices:}
\begin{enumerate}
\item \textbf{Structures de Contrôle}
   \begin{itemize}
   \item Dangling else et ambiguïté
   \item Shadowing et portée de blocs
   \item Fall-through dans switch
   \item Goto légitime (gestion d'erreurs)
   \end{itemize}

\item \textbf{Priorité et Conversions}
   \begin{itemize}
   \item Évaluer expressions complexes
   \item Piège signed/unsigned
   \item Promotion entière
   \item Conversions arithmétiques usuelles
   \end{itemize}

\item \textbf{Points de Séquence et UB}
   \begin{itemize}
   \item Modifications multiples entre points de séquence
   \item Ordre d'évaluation des arguments
   \item Débordement d'entiers (signed vs unsigned)
   \end{itemize}

\item \textbf{Déclaration/Définition}
   \begin{itemize}
   \item One Definition Rule (ODR)
   \item Static vs extern
   \end{itemize}

\item \textbf{Classes de Stockage}
   \begin{itemize}
   \item Static local (conservation entre appels)
   \item Const et volatile
   \end{itemize}

\item \textbf{Court-circuit}
   \begin{itemize}
   \item Sécurité avec \&\& et ||
   \item Division par zéro évitée
   \end{itemize}
\end{enumerate}

\textbf{TD\_Corrigé:} Solutions détaillées avec explications

\subsection*{TP: Travaux Pratiques (4-5 pages)}

\textbf{Environnement:} Docker (TP-1 complété)

\textbf{Exercices pratiques:}
\begin{enumerate}
\item \textbf{Phases de Compilation}
   \begin{itemize}
   \item Preprocessing: gcc -E (examiner .i)
   \item Compilation: gcc -S (examiner .s assembleur)
   \item Assemblage: gcc -c (générer .o)
   \item Linking: gcc (créer exécutable)
   \item Analyse avec nm et file
   \end{itemize}

\item \textbf{Options de Compilation}
   \begin{itemize}
   \item Compiler avec/sans warnings (-Wall -Wextra)
   \item Niveaux d'optimisation (-O0, -O2, -O3)
   \item Comparer tailles et performances
   \end{itemize}

\item \textbf{Debugging avec GDB}
   \begin{itemize}
   \item Compilation avec -g
   \item Breakpoints, step, next
   \item Inspection de variables (print)
   \item Pile d'appels (backtrace)
   \item Exemple avec fonction récursive
   \end{itemize}

\item \textbf{Projet Simple}
   \begin{itemize}
   \item Calculatrice avec opérations basiques
   \item Utilise seulement: int, float, boucles, fonctions
   \item PAS de tableaux, PAS de fichiers, PAS de malloc
   \end{itemize}
\end{enumerate}

\textbf{TP\_Pistes:} Indices pour le projet

\textbf{Durée estimée:} 3-4 heures

\newpage

\begin{partbox}
PARTIE II: REPRÉSENTATION DES DONNÉES
\end{partbox}

\section{Ch1: Représentation Numérique}

\textbf{Référence:} Bryant \& O'Hallaron - Chapitre 2

\subsection*{CM: Cours Magistral (10-12 pages)}

\textbf{Contenu théorique:}
\begin{enumerate}
\item \textbf{Systèmes de Numération}
   \begin{itemize}
   \item Binaire (base 2), octale (base 8), hexadécimale (base 16)
   \item Conversions entre bases
   \item Notation hexadécimale: TOUJOURS 2 caractères par octet (0A, FF, 00)
   \end{itemize}

\item \textbf{Représentation des Entiers Non-Signés}
   \begin{itemize}
   \item Encodage binaire pur
   \item Plage: $[0,\ 2^n - 1]$
   \item Arithmétique modulo $2^n$
   \end{itemize}

\item \textbf{Représentation des Entiers Signés (Complément à Deux)}
   \begin{itemize}
   \item Principe du complément à deux
   \item Plage: $[-2^{n-1},\ 2^{n-1} - 1]$
   \item Bit de signe
   \item Opérations: négation, addition, soustraction
   \end{itemize}

\item \textbf{Conversions Signed $\leftrightarrow$ Unsigned}
   \begin{itemize}
   \item Interprétation du pattern de bits
   \item Pièges et dangers
   \item Exemples de bugs critiques
   \end{itemize}

\item \textbf{Extension de Signe}
   \begin{itemize}
   \item Extension de 8 bits à 16 bits, 16 à 32, etc.
   \item Signed vs unsigned extension
   \end{itemize}

\item \textbf{Représentation IEEE 754 (Flottants)}
   \begin{itemize}
   \item Format: signe, exposant, mantisse
   \item float (32 bits) et double (64 bits)
   \item Valeurs spéciales: +0, -0, +$\infty$, -$\infty$, NaN
   \item Précision et erreurs d'arrondi
   \item Pourquoi ne jamais comparer avec ==
   \end{itemize}

\item \textbf{Dépassements et Débordements}
   \begin{itemize}
   \item Overflow/underflow signés (UB)
   \item Overflow/underflow non-signés (défini)
   \item Détection et prévention
   \end{itemize}
\end{enumerate}

\textbf{Visualisations TikZ (5-8 minimum):}
\begin{itemize}
\item Conversions binaire $\leftrightarrow$ hex $\leftrightarrow$ décimal
\item Représentation mémoire de valeurs signées/non-signées
\item Structure IEEE 754 (décomposition bits)
\item Exemples de conversions signed/unsigned avec patterns de bits
\item Dumps hexadécimaux avec 2 chars par octet
\end{itemize}

\subsection*{TD: Travaux Dirigés (4-5 pages)}

\textbf{Focus: Exercices forensics avec dumps hexadécimaux}

\textbf{Exercices:}
\begin{enumerate}
\item \textbf{Conversions de Bases}
   \begin{itemize}
   \item Binaire $\rightarrow$ Hex $\rightarrow$ Décimal
   \item Hex $\rightarrow$ Binaire $\rightarrow$ Décimal (signed/unsigned)
   \end{itemize}

\item \textbf{Complément à Deux}
   \begin{itemize}
   \item Calculer la négation
   \item Identifier la plage pour n bits
   \item Détecter overflow
   \end{itemize}

\item \textbf{Forensics Mémoire}
   \begin{itemize}
   \item Dumps hex: interpréter comme signed/unsigned
   \item Prédire les valeurs en mémoire
   \item Exemples: 0xFF = -1 (signed) ou 255 (unsigned)
   \end{itemize}

\item \textbf{IEEE 754}
   \begin{itemize}
   \item Décomposer float en bits (signe, exposant, mantisse)
   \item Calculer la valeur décimale
   \item Identifier NaN, infinités
   \end{itemize}

\item \textbf{Pièges Conversions}
   \begin{itemize}
   \item Prévoir le résultat de conversions signed/unsigned
   \item Bugs classiques (comparaisons, boucles)
   \end{itemize}
\end{enumerate}

\textbf{TD\_Corrigé:} Solutions avec dumps hex complets (2 chars/octet)

\subsection*{TP: Travaux Pratiques (5-6 pages)}

\textbf{Exercices pratiques:}
\begin{enumerate}
\item \textbf{Conversions et Visualisation}
   \begin{itemize}
   \item Programmes pour convertir dec/hex/bin
   \item Afficher représentation mémoire avec dump hex
   \end{itemize}

\item \textbf{Exploration Signed/Unsigned}
   \begin{itemize}
   \item Comparer comportements
   \item Reproduire bugs classiques
   \end{itemize}

\item \textbf{IEEE 754 Exploration}
   \begin{itemize}
   \item Décomposer float/double en bits
   \item Tester valeurs spéciales (NaN, infinités)
   \end{itemize}

\item \textbf{Projet: Bibliothèque de Conversions}
   \begin{itemize}
   \item Implémenter convertisseur universel
   \item Détection d'overflow
   \item Tests exhaustifs
   \end{itemize}
\end{enumerate}

\textbf{TP\_Pistes:} Indices pour le projet

\newpage

\begin{partbox}
PARTIE III: MÉMOIRE, POINTEURS, TABLEAUX
\end{partbox}

\section{Ch2: Organisation Mémoire}

\textbf{Référence:} Bryant Ch.2 (endianness) + minimal assembly

\subsection*{CM: Cours Magistral (10-12 pages)}

\textbf{Contenu théorique:}
\begin{enumerate}
\item \textbf{Modèle Mémoire Linéaire}
   \begin{itemize}
   \item Adresses et octets
   \item Espace d'adressage 32-bit vs 64-bit
   \item Segments: text, data, bss, heap, stack
   \end{itemize}

\item \textbf{Endianness}
   \begin{itemize}
   \item Little-endian (x86, x86-64)
   \item Big-endian (anciennes architectures RISC)
   \item Détection programmatique
   \item Impact sur le code portable
   \end{itemize}

\item \textbf{Adressage et Tailles}
   \begin{itemize}
   \item sizeof() pour types de base
   \item Architecture-dependent vs portable types
   \item Adresses: format hexadécimal
   \end{itemize}

\item \textbf{Introduction Minimale à l'Assembleur}
   \begin{itemize}
   \item Registres de base (rax, rbx, rsp, rbp)
   \item Instructions essentielles: mov, add, sub, lea
   \item Suffixes: b (byte), w (word), l (long), q (quad)
   \item Correspondance C $\rightarrow$ assembleur (exemples simples)
   \item \textbf{Note:} PAS un cours complet d'assembleur
   \end{itemize}

\item \textbf{Représentation en Mémoire}
   \begin{itemize}
   \item Variables globales vs locales
   \item Ordre en mémoire
   \item Dumps hexadécimaux avec GDB
   \end{itemize}
\end{enumerate}

\textbf{Visualisations TikZ (6-8):}
\begin{itemize}
\item Modèle mémoire linéaire (segments)
\item Endianness: disposition de 0x12345678 en little/big endian
\item Correspondance C/assembleur (exemples mov, add)
\item Dumps mémoire hexadécimaux annotés
\end{itemize}

\subsection*{TD: Travaux Dirigés (4-5 pages)}

\textbf{Exercices:}
\begin{enumerate}
\item \textbf{Endianness}
   \begin{itemize}
   \item Prédire disposition mémoire little/big endian
   \item Écrire détecteur d'endianness
   \end{itemize}

\item \textbf{Forensics Mémoire}
   \begin{itemize}
   \item Dumps hex: reconstruire valeurs multi-octets
   \item Identifier endianness à partir de dump
   \end{itemize}

\item \textbf{Assembleur Minimal}
   \begin{itemize}
   \item Lire code assembleur simple (mov, add)
   \item Correspondre avec code C
   \end{itemize}

\item \textbf{Sizeof et Adresses}
   \begin{itemize}
   \item Calculer tailles mémoire
   \item Prédire adresses relatives
   \end{itemize}
\end{enumerate}

\textbf{TD\_Corrigé:} Solutions avec dumps hex

\subsection*{TP: Travaux Pratiques (5-6 pages)}

\textbf{Exercices:}
\begin{enumerate}
\item \textbf{Exploration Endianness}
   \begin{itemize}
   \item Détecteur d'endianness
   \item Visualiser disposition mémoire
   \end{itemize}

\item \textbf{GDB et Dumps Mémoire}
   \begin{itemize}
   \item Commande x/Nx (examiner mémoire)
   \item Analyser variables en hex
   \end{itemize}

\item \textbf{Code C $\rightarrow$ Assembleur}
   \begin{itemize}
   \item Compiler avec gcc -S
   \item Analyser .s généré
   \item Identifier correspondances
   \end{itemize}

\item \textbf{Projet: Memory Inspector}
   \begin{itemize}
   \item Afficher mémoire en hex
   \item Interpréter comme différents types
   \end{itemize}
\end{enumerate}

\textbf{TP\_Pistes:} Indices projet

\newpage

\section{Ch3: Pointeurs}

\textbf{Référence:} Bryant Ch.3 (pointer arithmetic)

\subsection*{CM: Cours Magistral (10-12 pages)}

\textbf{Contenu théorique:}
\begin{enumerate}
\item \textbf{Concepts de Base}
   \begin{itemize}
   \item Définition: adresse mémoire
   \item Déclaration: int *p;
   \item Opérateurs: \& (adresse) et * (déréférence)
   \item NULL et pointeur invalide
   \end{itemize}

\item \textbf{Arithmétique de Pointeurs}
   \begin{itemize}
   \item p + n : avance de n * sizeof(*p) octets
   \item p - q : nombre d'éléments entre deux pointeurs
   \item Incrémentation/décrémentation
   \item Comparaison de pointeurs
   \end{itemize}

\item \textbf{Pointeurs et Fonctions}
   \begin{itemize}
   \item Passage par référence simulé
   \item Retour de pointeurs (dangers)
   \item Pointeurs vers fonctions (aperçu)
   \end{itemize}

\item \textbf{Pointeurs de Pointeurs}
   \begin{itemize}
   \item Déclaration: int **pp;
   \item Usage et utilité
   \end{itemize}

\item \textbf{Pointeurs const}
   \begin{itemize}
   \item const int *p vs int * const p
   \item const int * const p
   \end{itemize}

\item \textbf{Pièges et Erreurs}
   \begin{itemize}
   \item Déréférence de NULL
   \item Pointeurs pendants (dangling)
   \item Pointeurs non initialisés
   \item Aliasing
   \end{itemize}
\end{enumerate}

\textbf{Visualisations TikZ (6-8):}
\begin{itemize}
\item Pointeur $\rightarrow$ case mémoire
\item Arithmétique: p, p+1, p+2 avec adresses
\item Pointeurs de pointeurs (double indirection)
\item Dumps hex avec pointeurs
\end{itemize}

\subsection*{TD: Travaux Dirigés (4-5 pages)}

\textbf{Exercices forensics:}
\begin{enumerate}
\item \textbf{Arithmétique de Pointeurs}
   \begin{itemize}
   \item Calculer p+n avec adresses
   \item Prédire sorties
   \end{itemize}

\item \textbf{Dumps Hexadécimaux}
   \begin{itemize}
   \item Reconstruire structure de pointeurs
   \item Suivre indirections
   \end{itemize}

\item \textbf{Erreurs Classiques}
   \begin{itemize}
   \item Identifier bugs (NULL, dangling)
   \item Corriger code défectueux
   \end{itemize}
\end{enumerate}

\textbf{TD\_Corrigé:} Solutions détaillées

\subsection*{TP: Travaux Pratiques (5-6 pages)}

\textbf{Exercices:}
\begin{enumerate}
\item \textbf{Manipulation de Base}
   \begin{itemize}
   \item Swap via pointeurs
   \item Fonctions modifiant via pointeurs
   \end{itemize}

\item \textbf{Arithmétique}
   \begin{itemize}
   \item Parcourir zone mémoire
   \item Calculer distances
   \end{itemize}

\item \textbf{Debugging}
   \begin{itemize}
   \item Traquer erreurs avec GDB
   \item Valgrind: invalid read/write
   \end{itemize}

\item \textbf{Projet: Generic Swap}
   \begin{itemize}
   \item Swap universel avec void*
   \item Tests sur différents types
   \end{itemize}
\end{enumerate}

\textbf{TP\_Pistes:} Indices

\newpage

\section{Ch4: Tableaux}

\textbf{Référence:} Bryant Ch.3 Section 3.8 + CSE 333

\subsection*{CM: Cours Magistral (10-12 pages)}

\textbf{Contenu théorique:}
\begin{enumerate}
\item \textbf{Tableaux 1D}
   \begin{itemize}
   \item Déclaration et initialisation
   \item Accès par indice: tab[i]
   \item Relation avec pointeurs: tab $\equiv$ \&tab[0]
   \item tab[i] $\equiv$ *(tab + i)
   \item Décomposition d'adresse
   \end{itemize}

\item \textbf{Passage de Tableaux aux Fonctions}
   \begin{itemize}
   \item Dégénérescence en pointeur
   \item Nécessité de passer la taille
   \item Modifications visibles dans l'appelant
   \end{itemize}

\item \textbf{Tableaux 2D}
   \begin{itemize}
   \item Déclaration: int mat[N][M];
   \item Organisation en mémoire (row-major)
   \item Adressage: mat[i][j]
   \item Arithmétique 2D: *(*(mat + i) + j)
   \end{itemize}

\item \textbf{Tableaux Multidimensionnels}
   \begin{itemize}
   \item Tableaux 3D et au-delà
   \item Calcul d'adresse
   \end{itemize}

\item \textbf{Pièges}
   \begin{itemize}
   \item Dépassement de bornes
   \item sizeof(tableau) vs sizeof(pointeur)
   \item Tableaux de taille variable (VLA)
   \end{itemize}
\end{enumerate}

\textbf{Visualisations TikZ (6-8):}
\begin{itemize}
\item Tableau 1D en mémoire avec adresses
\item Relation tab[i] $\equiv$ *(tab+i)
\item Tableau 2D: row-major layout
\item Dumps hex de tableaux
\end{itemize}

\subsection*{TD: Travaux Dirigés (4-5 pages)}

\textbf{Exercices forensics:}
\begin{enumerate}
\item \textbf{Arithmétique Tableaux}
   \begin{itemize}
   \item Prédire adresses
   \item Calculs avec pointeurs
   \end{itemize}

\item \textbf{Tableaux 2D}
   \begin{itemize}
   \item Layout mémoire
   \item Accès par arithmétique
   \end{itemize}

\item \textbf{Dumps Hexadécimaux}
   \begin{itemize}
   \item Reconstruire tableau depuis dump
   \item Identifier dimensions
   \end{itemize}

\item \textbf{Bugs Classiques}
   \begin{itemize}
   \item Dépassement de bornes
   \item sizeof confusion
   \end{itemize}
\end{enumerate}

\textbf{TD\_Corrigé:} Solutions avec dumps

\subsection*{TP: Travaux Pratiques (5-6 pages)}

\textbf{Exercices:}
\begin{enumerate}
\item \textbf{Manipulation 1D}
   \begin{itemize}
   \item Recherche, tri, statistiques
   \item Deux méthodes: indices et pointeurs
   \end{itemize}

\item \textbf{Tableaux 2D}
   \begin{itemize}
   \item Opérations matricielles simples
   \item Transposition, addition
   \end{itemize}

\item \textbf{Projet: Image Processing}
   \begin{itemize}
   \item Tableau 2D comme image
   \item Filtres simples (moyenne, seuillage)
   \end{itemize}
\end{enumerate}

\textbf{TP\_Pistes:} Indices

\newpage

\section{Ch5: Chaînes de Caractères}

\textbf{Référence:} Bryant Ch.3 + CSE 333

\subsection*{CM: Cours Magistral (10-12 pages)}

\textbf{Contenu théorique:}
\begin{enumerate}
\item \textbf{Représentation}
   \begin{itemize}
   \item Tableau de char terminé par '\textbackslash 0'
   \item char str[] vs char *str
   \item Chaînes littérales (constantes)
   \end{itemize}

\item \textbf{Fonctions Standard}
   \begin{itemize}
   \item strlen, strcpy, strncpy
   \item strcmp, strncmp
   \item strcat, strncat
   \item strchr, strstr
   \end{itemize}

\item \textbf{Manipulation Manuelle}
   \begin{itemize}
   \item Parcours avec pointeurs
   \item Modification caractère par caractère
   \end{itemize}

\item \textbf{Sécurité}
   \begin{itemize}
   \item Buffer overflow
   \item Fonctions dangereuses: strcpy, strcat, gets
   \item Fonctions sûres: strncpy, strncat, fgets
   \item Validation d'entrée
   \end{itemize}

\item \textbf{Conversions}
   \begin{itemize}
   \item atoi, atof, strtol, strtod
   \item sprintf, snprintf
   \end{itemize}
\end{enumerate}

\textbf{Visualisations TikZ (5-7):}
\begin{itemize}
\item Chaîne en mémoire avec '\textbackslash 0'
\item Buffer overflow (avant/après)
\item Dumps hex de chaînes
\end{itemize}

\subsection*{TD: Travaux Dirigés (4-5 pages)}

\textbf{Exercices:}
\begin{enumerate}
\item \textbf{Manipulation}
   \begin{itemize}
   \item Implémenter strlen, strcpy manuellement
   \item Prédire sorties
   \end{itemize}

\item \textbf{Sécurité}
   \begin{itemize}
   \item Identifier buffer overflows
   \item Corriger code vulnérable
   \end{itemize}

\item \textbf{Forensics}
   \begin{itemize}
   \item Dumps hex de chaînes
   \item Reconstruire texte
   \end{itemize}
\end{enumerate}

\textbf{TD\_Corrigé:} Solutions

\subsection*{TP: Travaux Pratiques (5-6 pages)}

\textbf{Exercices:}
\begin{enumerate}
\item \textbf{Bibliothèque Chaînes}
   \begin{itemize}
   \item Réimplémenter fonctions standard
   \item Tests exhaustifs
   \end{itemize}

\item \textbf{Parsing}
   \begin{itemize}
   \item Tokenisation
   \item Extraction de données
   \end{itemize}

\item \textbf{Projet: Text Editor Simple}
   \begin{itemize}
   \item Opérations sur texte
   \item Sécurisé (pas de overflow)
   \end{itemize}
\end{enumerate}

\textbf{TP\_Pistes:} Indices

\newpage

\begin{partbox}
PARTIE IV: STRUCTURES ET ALLOCATION
\end{partbox}

\section{Ch6: Structures et Unions}

\textbf{Référence:} Bryant Ch.3 Section 3.9

\subsection*{CM: Cours Magistral (10-12 pages)}

\textbf{Contenu théorique:}
\begin{enumerate}
\item \textbf{Structures}
   \begin{itemize}
   \item Déclaration, définition
   \item Initialisation
   \item Accès membres: . et ->
   \item typedef
   \end{itemize}

\item \textbf{Alignment et Padding}
   \begin{itemize}
   \item Règles d'alignement (architecture-dependent)
   \item Calcul de padding
   \item sizeof(struct) vs somme des membres
   \item Impact sur performances
   \end{itemize}

\item \textbf{Structures Imbriquées}
   \begin{itemize}
   \item Struct dans struct
   \item Tableaux de structures
   \item Pointeurs vers structures
   \end{itemize}

\item \textbf{Unions}
   \begin{itemize}
   \item Définition et usage
   \item Tous les membres partagent la mémoire
   \item Cas d'usage
   \end{itemize}

\item \textbf{Bit Fields}
   \begin{itemize}
   \item Déclaration
   \item Usage (registres hardware, protocoles)
   \item Portabilité
   \end{itemize}

\item \textbf{Structures Opaques}
   \begin{itemize}
   \item Forward declaration
   \item Encapsulation
   \end{itemize}
\end{enumerate}

\textbf{Visualisations TikZ (6-8):}
\begin{itemize}
\item Structure avec padding (annotation octets)
\item Dumps hex de structures
\item Union: overlap mémoire
\item Alignement: exemples avant/après réorganisation
\end{itemize}

\subsection*{TD: Travaux Dirigés (4-5 pages)}

\textbf{Exercices forensics:}
\begin{enumerate}
\item \textbf{Calcul Padding}
   \begin{itemize}
   \item Prédire sizeof(struct)
   \item Identifier padding
   \end{itemize}

\item \textbf{Dumps Hexadécimaux}
   \begin{itemize}
   \item Reconstruire struct depuis dump
   \item Identifier alignment
   \end{itemize}

\item \textbf{Optimisation}
   \begin{itemize}
   \item Réorganiser membres pour réduire padding
   \end{itemize}

\item \textbf{Unions}
   \begin{itemize}
   \item Comprendre overlap
   \item Cas d'usage
   \end{itemize}
\end{enumerate}

\textbf{TD\_Corrigé:} Solutions

\subsection*{TP: Travaux Pratiques (5-6 pages)}

\textbf{Exercices:}
\begin{enumerate}
\item \textbf{Manipulation Structures}
   \begin{itemize}
   \item Créer, initialiser, copier
   \item Tableaux de structures
   \end{itemize}

\item \textbf{Analyse Mémoire}
   \begin{itemize}
   \item GDB: examiner structures
   \item Dumps hex
   \end{itemize}

\item \textbf{Projet: Database Simple}
   \begin{itemize}
   \item Gestion d'enregistrements
   \item Recherche, tri
   \end{itemize}
\end{enumerate}

\textbf{TP\_Pistes:} Indices

\newpage

\section{Ch7: Allocation Dynamique}

\textbf{Référence:} Bryant Ch.9 (Virtual Memory - malloc/free)

\subsection*{CM: Cours Magistral (10-12 pages)}

\textbf{Contenu théorique:}
\begin{enumerate}
\item \textbf{Heap et Allocation}
   \begin{itemize}
   \item Différence stack vs heap
   \item malloc, calloc, realloc, free
   \item Durée de vie des allocations
   \end{itemize}

\item \textbf{Gestion de la Mémoire}
   \begin{itemize}
   \item Fragmentation interne/externe
   \item Stratégies d'allocation (first-fit, best-fit)
   \item Coalescing
   \end{itemize}

\item \textbf{Erreurs Classiques}
   \begin{itemize}
   \item Memory leaks
   \item Double free
   \item Use after free
   \item Buffer overflow
   \item Pointeurs pendants
   \end{itemize}

\item \textbf{Valgrind}
   \begin{itemize}
   \item Détection de leaks
   \item Invalid read/write
   \item Utilisation mémoire non initialisée
   \end{itemize}

\item \textbf{Allocations Complexes}
   \begin{itemize}
   \item Tableaux dynamiques 1D
   \item Tableaux dynamiques 2D (deux méthodes)
   \item Structures dynamiques
   \end{itemize}
\end{enumerate}

\textbf{Visualisations TikZ (6-8):}
\begin{itemize}
\item Heap vs stack
\item Fragmentation (schémas)
\item Erreurs: leak, double free, use after free
\item Dumps hex avec metadata malloc
\end{itemize}

\subsection*{TD: Travaux Dirigés (4-5 pages)}

\textbf{Exercices:}
\begin{enumerate}
\item \textbf{Allocation/Libération}
   \begin{itemize}
   \item Identifier leaks
   \item Corriger erreurs
   \end{itemize}

\item \textbf{Tableaux Dynamiques}
   \begin{itemize}
   \item Allocation 2D
   \item Libération correcte
   \end{itemize}

\item \textbf{Forensics Valgrind}
   \begin{itemize}
   \item Interpréter rapports
   \item Localiser bugs
   \end{itemize}
\end{enumerate}

\textbf{TD\_Corrigé:} Solutions

\subsection*{TP: Travaux Pratiques (5-6 pages)}

\textbf{Exercices:}
\begin{enumerate}
\item \textbf{Détection Erreurs}
   \begin{itemize}
   \item Programmes bugués
   \item Valgrind pour diagnostiquer
   \end{itemize}

\item \textbf{Structures Dynamiques}
   \begin{itemize}
   \item Listes chaînées
   \item Arbres simples
   \end{itemize}

\item \textbf{Projet: Dynamic Array Library}
   \begin{itemize}
   \item Array redimensionnable
   \item Tests Valgrind (0 leaks)
   \end{itemize}
\end{enumerate}

\textbf{TP\_Pistes:} Indices

\newpage

\begin{partbox}
PARTIE V: EXÉCUTION ET SYSTÈME
\end{partbox}

\section{Ch8: Stack et Calling Conventions}

\textbf{Référence:} Bryant Ch.3 (Machine-Level Programming)

\subsection*{CM: Cours Magistral (10-12 pages)}

\textbf{Contenu théorique:}
\begin{enumerate}
\item \textbf{Organisation de la Pile}
   \begin{itemize}
   \item Croissance vers adresses basses
   \item Stack pointer (rsp)
   \item Base pointer (rbp)
   \end{itemize}

\item \textbf{Stack Frames}
   \begin{itemize}
   \item Structure: paramètres, adresse retour, variables locales
   \item Prologue et épilogue de fonction
   \item Frame pointer (optionnel avec -O2)
   \end{itemize}

\item \textbf{Calling Conventions (x86-64 System V ABI)}
   \begin{itemize}
   \item Passage paramètres: registres (rdi, rsi, rdx, rcx, r8, r9) puis stack
   \item Valeur de retour: rax
   \item Registres caller-saved vs callee-saved
   \end{itemize}

\item \textbf{Récursivité}
   \begin{itemize}
   \item Stack frames multiples
   \item Limites et stack overflow
   \end{itemize}

\item \textbf{Buffer Overflow sur la Stack}
   \begin{itemize}
   \item Écrasement adresse de retour
   \item Protections: stack canaries, ASLR
   \end{itemize}
\end{enumerate}

\textbf{Visualisations TikZ (6-8):}
\begin{itemize}
\item Stack frame (détaillé)
\item Évolution stack lors d'appels
\item Récursivité: frames empilés
\item Buffer overflow: avant/après
\end{itemize}

\subsection*{TD: Travaux Dirigés (4-5 pages)}

\textbf{Exercices forensics:}
\begin{enumerate}
\item \textbf{Stack Frames}
   \begin{itemize}
   \item Dessiner état stack
   \item Tracer appels de fonction
   \end{itemize}

\item \textbf{Calling Conventions}
   \begin{itemize}
   \item Identifier registres utilisés
   \item Paramètres: registres vs stack
   \end{itemize}

\item \textbf{Dumps GDB}
   \begin{itemize}
   \item Examiner stack avec GDB
   \item Interpréter backtrace
   \end{itemize}
\end{enumerate}

\textbf{TD\_Corrigé:} Solutions

\subsection*{TP: Travaux Pratiques (5-6 pages)}

\textbf{Exercices:}
\begin{enumerate}
\item \textbf{Exploration avec GDB}
   \begin{itemize}
   \item Examiner frames (info frame)
   \item Remonter la pile (up/down)
   \end{itemize}

\item \textbf{Assembleur}
   \begin{itemize}
   \item Compiler avec -S
   \item Identifier prologue/épilogue
   \end{itemize}

\item \textbf{Projet: Stack Tracer}
   \begin{itemize}
   \item Afficher call stack
   \item Analyser récursivité
   \end{itemize}
\end{enumerate}

\textbf{TP\_Pistes:} Indices

\newpage

\section{Ch9: Memory Layout et Cache}

\textbf{Référence:} Bryant Ch.6 (Memory Hierarchy) + Ch.9 (Virtual Memory)

\subsection*{CM: Cours Magistral (10-12 pages)}

\textbf{Contenu théorique:}
\begin{enumerate}
\item \textbf{Segments Mémoire}
   \begin{itemize}
   \item Text: code exécutable
   \item Data: variables globales initialisées
   \item BSS: variables globales non initialisées
   \item Heap: allocation dynamique
   \item Stack: variables locales, appels
   \end{itemize}

\item \textbf{Hiérarchie Mémoire}
   \begin{itemize}
   \item Registres, L1, L2, L3, RAM, disque
   \item Latences et bande passante
   \item Principe de localité (temporelle, spatiale)
   \end{itemize}

\item \textbf{Cache}
   \begin{itemize}
   \item Cache lines (typiquement 64 octets)
   \item Direct-mapped, set-associative, fully-associative
   \item Cache hits vs misses
   \item Cold, capacity, conflict misses
   \end{itemize}

\item \textbf{Optimisation Cache}
   \begin{itemize}
   \item Accès séquentiel vs random
   \item Blocking/tiling pour matrices
   \item Cache-friendly data structures
   \end{itemize}

\item \textbf{Mémoire Virtuelle (introduction)}
   \begin{itemize}
   \item Espace d'adressage virtuel vs physique
   \item Pages et page tables
   \item TLB (Translation Lookaside Buffer)
   \end{itemize}
\end{enumerate}

\textbf{Visualisations TikZ (6-8):}
\begin{itemize}
\item Layout mémoire complet (segments)
\item Hiérarchie mémoire (pyramide)
\item Cache: structure et fonctionnement
\item Accès matrice: row-major vs column-major
\end{itemize}

\subsection*{TD: Travaux Dirigés (4-5 pages)}

\textbf{Exercices:}
\begin{enumerate}
\item \textbf{Segments}
   \begin{itemize}
   \item Identifier où variables sont stockées
   \item Calculer tailles segments
   \end{itemize}

\item \textbf{Cache}
   \begin{itemize}
   \item Calculer hits/misses pour accès donnés
   \item Analyser patterns d'accès
   \end{itemize}

\item \textbf{Optimisation}
   \begin{itemize}
   \item Comparer versions (cache-friendly vs non)
   \end{itemize}
\end{enumerate}

\textbf{TD\_Corrigé:} Solutions

\subsection*{TP: Travaux Pratiques (5-6 pages)}

\textbf{Exercices:}
\begin{enumerate}
\item \textbf{Mesure Performances}
   \begin{itemize}
   \item Accès séquentiel vs random
   \item Impact taille cache
   \end{itemize}

\item \textbf{Optimisation Matrices}
   \begin{itemize}
   \item Transposition cache-aware
   \item Multiplication avec blocking
   \end{itemize}

\item \textbf{Projet: Cache Simulator}
   \begin{itemize}
   \item Simuler cache simple
   \item Mesurer hits/misses
   \end{itemize}
\end{enumerate}

\textbf{TP\_Pistes:} Indices

\newpage

\section{Ch10: POSIX File I/O}

\textbf{Référence:} CSE 333 (System calls)

\subsection*{CM: Cours Magistral (10-12 pages)}

\textbf{Contenu théorique:}
\begin{enumerate}
\item \textbf{File Descriptors}
   \begin{itemize}
   \item Concept: entiers identifiant fichiers ouverts
   \item stdin (0), stdout (1), stderr (2)
   \item Limite par processus
   \end{itemize}

\item \textbf{System Calls}
   \begin{itemize}
   \item open, close
   \item read, write
   \item lseek (positionnement)
   \item Flags: O\_RDONLY, O\_WRONLY, O\_RDWR, O\_CREAT, O\_APPEND, etc.
   \item Permissions (mode)
   \end{itemize}

\item \textbf{Unbuffered vs Buffered I/O}
   \begin{itemize}
   \item System calls: accès direct kernel
   \item stdio (FILE*): bufferisé
   \item Performances et cas d'usage
   \end{itemize}

\item \textbf{Gestion d'Erreurs}
   \begin{itemize}
   \item Codes de retour (-1 pour erreurs)
   \item errno et perror
   \item strerror
   \end{itemize}

\item \textbf{Opérations Avancées}
   \begin{itemize}
   \item dup, dup2 (duplication descripteurs)
   \item fcntl (contrôle)
   \item Redirections
   \end{itemize}

\item \textbf{Directories}
   \begin{itemize}
   \item opendir, readdir, closedir
   \item stat, fstat
   \end{itemize}
\end{enumerate}

\textbf{Visualisations TikZ (5-7):}
\begin{itemize}
\item Table de descripteurs par processus
\item Buffering: unbuffered vs buffered
\item Redirections avec dup2
\end{itemize}

\subsection*{TD: Travaux Dirigés (4-5 pages)}

\textbf{Exercices:}
\begin{enumerate}
\item \textbf{System Calls}
   \begin{itemize}
   \item Écrire programmes simples (copie fichier)
   \item Gestion d'erreurs
   \end{itemize}

\item \textbf{File Descriptors}
   \begin{itemize}
   \item Tracer état descripteurs
   \item Comprendre redirections
   \end{itemize}

\item \textbf{Performances}
   \begin{itemize}
   \item Comparer read/write vs fread/fwrite
   \end{itemize}
\end{enumerate}

\textbf{TD\_Corrigé:} Solutions

\subsection*{TP: Travaux Pratiques (5-6 pages)}

\textbf{Exercices:}
\begin{enumerate}
\item \textbf{Manipulation Fichiers}
   \begin{itemize}
   \item Copie, concaténation
   \item Statistiques (taille, permissions)
   \end{itemize}

\item \textbf{Directories}
   \begin{itemize}
   \item Lister récursivement
   \item Filtrer par extension
   \end{itemize}

\item \textbf{Projet: File Manager}
   \begin{itemize}
   \item Mini gestionnaire de fichiers
   \item Opérations POSIX
   \end{itemize}
\end{enumerate}

\textbf{TP\_Pistes:} Indices

\newpage

\begin{partbox}
PARTIE VI: OUTILS ET TECHNIQUES AVANCÉES
\end{partbox}

\section{Ch11: Makefiles et Linking}

\textbf{Référence:} CSE 333 (Build systems) + Bryant Ch.7 (Linking)

\subsection*{CM: Cours Magistral (10-12 pages)}

\textbf{Contenu théorique:}
\begin{enumerate}
\item \textbf{Makefiles}
   \begin{itemize}
   \item Syntaxe: cibles, dépendances, commandes
   \item Variables ($CC, $CFLAGS, etc.)
   \item Règles implicites et patterns
   \item Phony targets (.PHONY)
   \item Makefile modulaire
   \end{itemize}

\item \textbf{Process de Linking}
   \begin{itemize}
   \item Symbol resolution
   \item Relocation
   \item Fichiers objets (.o)
   \item Archives statiques (.a)
   \end{itemize}

\item \textbf{Bibliothèques Statiques}
   \begin{itemize}
   \item Création avec ar
   \item Linking avec -l et -L
   \item Ordre de linking
   \end{itemize}

\item \textbf{Bibliothèques Dynamiques (Shared Libraries)}
   \begin{itemize}
   \item Création (.so sous Linux, .dll sous Windows)
   \item Position-Independent Code (-fPIC)
   \item Loading dynamique (dlopen, dlsym)
   \item LD\_LIBRARY\_PATH
   \end{itemize}

\item \textbf{Outils}
   \begin{itemize}
   \item nm, objdump, ldd
   \item Analyse de symboles et dépendances
   \end{itemize}
\end{enumerate}

\textbf{Visualisations TikZ (5-7):}
\begin{itemize}
\item Graphe de dépendances Makefile
\item Process de linking (schéma)
\item Bibliothèques statiques vs dynamiques
\end{itemize}

\subsection*{TD: Travaux Dirigés (3-4 pages)}

\textbf{Exercices:}
\begin{enumerate}
\item \textbf{Makefiles}
   \begin{itemize}
   \item Écrire Makefile pour projet
   \item Optimiser avec règles patterns
   \end{itemize}

\item \textbf{Linking}
   \begin{itemize}
   \item Résoudre erreurs de linking
   \item Ordre correct bibliothèques
   \end{itemize}

\item \textbf{Bibliothèques}
   \begin{itemize}
   \item Créer .a et .so
   \item Lier programmes avec
   \end{itemize}
\end{enumerate}

\textbf{TD\_Corrigé:} Solutions

\subsection*{TP: Travaux Pratiques (5-6 pages)}

\textbf{Exercices:}
\begin{enumerate}
\item \textbf{Makefile Complexe}
   \begin{itemize}
   \item Projet multi-fichiers
   \item Règles de nettoyage
   \end{itemize}

\item \textbf{Bibliothèques}
   \begin{itemize}
   \item Créer bibliothèque statique
   \item Créer bibliothèque dynamique
   \item Programmes clients
   \end{itemize}

\item \textbf{Projet: Build System}
   \begin{itemize}
   \item Système de build complet
   \item Tests, documentation
   \end{itemize}
\end{enumerate}

\textbf{TP\_Pistes:} Indices

\newpage

\section{Ch12: Function Pointers}

\textbf{Référence:} Bryant Ch.3 + patterns

\subsection*{CM: Cours Magistral (8-10 pages)}

\textbf{Contenu théorique:}
\begin{enumerate}
\item \textbf{Concept}
   \begin{itemize}
   \item Pointeurs vers fonctions
   \item Déclaration: int (*func\_ptr)(int, int);
   \item Assignation et appel
   \end{itemize}

\item \textbf{Typedef pour Lisibilité}
   \begin{itemize}
   \item typedef int (*operation)(int, int);
   \item Simplification code
   \end{itemize}

\item \textbf{Callbacks}
   \begin{itemize}
   \item Passer fonction en paramètre
   \item Exemples: qsort, bsearch
   \end{itemize}

\item \textbf{Dispatch Tables}
   \begin{itemize}
   \item Tableaux de pointeurs de fonctions
   \item Switch alternatif
   \item State machines
   \end{itemize}

\item \textbf{Applications}
   \begin{itemize}
   \item Plugins et extensibilité
   \item Event handlers
   \item Algorithmes génériques
   \end{itemize}
\end{enumerate}

\textbf{Visualisations TikZ (4-6):}
\begin{itemize}
\item Pointeur de fonction $\rightarrow$ code
\item Dispatch table (schéma)
\item Callback (flux d'exécution)
\end{itemize}

\subsection*{TD: Travaux Dirigés (3-4 pages)}

\textbf{Exercices:}
\begin{enumerate}
\item \textbf{Déclarations}
   \begin{itemize}
   \item Déchiffrer déclarations complexes
   \item Écrire déclarations
   \end{itemize}

\item \textbf{Callbacks}
   \begin{itemize}
   \item Utiliser qsort avec comparateurs
   \item Écrire comparateurs personnalisés
   \end{itemize}

\item \textbf{Dispatch Tables}
   \begin{itemize}
   \item Implémenter calculatrice
   \item State machine simple
   \end{itemize}
\end{enumerate}

\textbf{TD\_Corrigé:} Solutions

\subsection*{TP: Travaux Pratiques (4-5 pages)}

\textbf{Exercices:}
\begin{enumerate}
\item \textbf{Algorithmes Génériques}
   \begin{itemize}
   \item Map, filter, reduce
   \item Tri avec callbacks
   \end{itemize}

\item \textbf{Event System}
   \begin{itemize}
   \item Enregistrement handlers
   \item Déclenchement événements
   \end{itemize}

\item \textbf{Projet: Plugin System}
   \begin{itemize}
   \item Chargement dynamique
   \item Interface via function pointers
   \end{itemize}
\end{enumerate}

\textbf{TP\_Pistes:} Indices

\newpage

\section{Ch13: Preprocessing et Optimisation}

\textbf{Référence:} Bryant Ch.5 (Optimizing Program Performance) + preprocessing

\subsection*{CM: Cours Magistral (10-12 pages)}

\textbf{Contenu théorique:}
\begin{enumerate}
\item \textbf{Preprocessing}
   \begin{itemize}
   \item Directives: \#include, \#define, \#ifdef, \#ifndef, \#endif
   \item Macros avec paramètres
   \item Pièges des macros (effets de bord, priorité)
   \item Macros vs inline functions
   \item Conditional compilation
   \item \#pragma
   \end{itemize}

\item \textbf{Include Guards}
   \begin{itemize}
   \item Prévenir inclusions multiples
   \item \#pragma once (non-standard mais répandu)
   \end{itemize}

\item \textbf{Optimisations Compilateur}
   \begin{itemize}
   \item Niveaux: -O0, -O1, -O2, -O3, -Os
   \item Optimisations courantes: inlining, loop unrolling, constant folding
   \item Profile-guided optimization (PGO)
   \end{itemize}

\item \textbf{Optimisations Manuel}
   \begin{itemize}
   \item Réduction appels de fonction
   \item Loop hoisting
   \item Strength reduction
   \item Cache-aware code
   \end{itemize}

\item \textbf{Inline Functions}
   \begin{itemize}
   \item Keyword inline
   \item Avantages vs macros
   \item Limitations
   \end{itemize}

\item \textbf{Profiling}
   \begin{itemize}
   \item gprof
   \item perf (Linux)
   \item Identifier hotspots
   \end{itemize}

\item \textbf{Trade-offs}
   \begin{itemize}
   \item Performance vs lisibilité
   \item Space vs time
   \item Premature optimization
   \end{itemize}
\end{enumerate}

\textbf{Visualisations TikZ (5-7):}
\begin{itemize}
\item Expansion de macros
\item Comparaison code non-optimisé vs optimisé
\item Impact cache (before/after optimization)
\end{itemize}

\subsection*{TD: Travaux Dirigés (4-5 pages)}

\textbf{Exercices:}
\begin{enumerate}
\item \textbf{Macros}
   \begin{itemize}
   \item Écrire macros correctes
   \item Identifier bugs dans macros
   \end{itemize}

\item \textbf{Preprocessing}
   \begin{itemize}
   \item Expansion manuelle
   \item Conditional compilation
   \end{itemize}

\item \textbf{Optimisations}
   \begin{itemize}
   \item Identifier opportunités
   \item Appliquer transformations
   \end{itemize}
\end{enumerate}

\textbf{TD\_Corrigé:} Solutions

\subsection*{TP: Travaux Pratiques (5-6 pages)}

\textbf{Exercices:}
\begin{enumerate}
\item \textbf{Analyse Preprocessing}
   \begin{itemize}
   \item Examiner .i files
   \item Comprendre expansions
   \end{itemize}

\item \textbf{Profiling}
   \begin{itemize}
   \item Utiliser gprof
   \item Identifier hotspots
   \end{itemize}

\item \textbf{Optimisation}
   \begin{itemize}
   \item Programme lent fourni
   \item Optimiser (manuel + compilateur)
   \item Mesurer gains
   \end{itemize}

\item \textbf{Projet: Performance Challenge}
   \begin{itemize}
   \item Implémenter algorithme
   \item Optimiser au maximum
   \item Benchmark
   \end{itemize}
\end{enumerate}

\textbf{TP\_Pistes:} Indices

\newpage

\begin{partbox}
CHAPITRES OPTIONNELS
\end{partbox}

\section{Ch14 (Optionnel): Multithreading Basics}

\textbf{Référence:} Bryant Ch.12 (Concurrent Programming) - introduction seulement

\subsection*{CM: Cours Magistral (8-10 pages)}

\textbf{Contenu théorique:}
\begin{itemize}
\item Threads vs processus
\item pthread\_create, pthread\_join
\item Problèmes: race conditions, deadlocks
\item Mutexes (pthread\_mutex\_t)
\item Introduction synchronisation
\end{itemize}

\textbf{Note:} Cours d'introduction uniquement, pas de traitement exhaustif.

\subsection*{TD \& TP: (3-4 pages chacun)}

Exercices simples de création de threads et synchronisation basique.

\newpage

\section{Ch15 (Optionnel): Networking Basics}

\textbf{Référence:} Bryant Ch.11 (Network Programming) - introduction seulement

\subsection*{CM: Cours Magistral (8-10 pages)}

\textbf{Contenu théorique:}
\begin{itemize}
\item Modèle client-serveur
\item Sockets: socket(), bind(), listen(), accept(), connect()
\item TCP vs UDP (aperçu)
\item Exemple simple: echo server
\end{itemize}

\textbf{Note:} Cours d'introduction uniquement, pas de traitement exhaustif.

\subsection*{TD \& TP: (3-4 pages chacun)}

Exercices simples de communication réseau.

\newpage

\section*{Résumé de la Structure du Cours}

\subsection*{Chapitres Obligatoires}

\begin{itemize}
\item \textbf{TP-1:} Setup environnement Docker + vim
\item \textbf{Ch-1:} Introduction (2-3 pages)
\item \textbf{Ch0:} Syntaxe C complète (fondamentaux)
\item \textbf{Ch1:} Représentation numérique
\item \textbf{Ch2:} Organisation mémoire
\item \textbf{Ch3:} Pointeurs
\item \textbf{Ch4:} Tableaux
\item \textbf{Ch5:} Chaînes
\item \textbf{Ch6:} Structures et unions
\item \textbf{Ch7:} Allocation dynamique
\item \textbf{Ch8:} Stack et calling conventions
\item \textbf{Ch9:} Memory layout et cache
\item \textbf{Ch10:} POSIX File I/O
\item \textbf{Ch11:} Makefiles et linking
\item \textbf{Ch12:} Function pointers
\item \textbf{Ch13:} Preprocessing et optimisation
\end{itemize}

\subsection*{Chapitres Optionnels}

\begin{itemize}
\item \textbf{Ch14:} Multithreading (introduction)
\item \textbf{Ch15:} Networking (introduction)
\end{itemize}

\subsection*{Volume Total}

\begin{itemize}
\item \textbf{CM:} 13 chapitres × 10 pages $\approx$ 130 pages
\item \textbf{TD:} 13 chapitres × 4 pages $\approx$ 52 pages
\item \textbf{TD\_Corrigé:} 13 chapitres × 4 pages $\approx$ 52 pages
\item \textbf{TP:} 13 chapitres × 5 pages $\approx$ 65 pages
\item \textbf{TP\_Pistes:} 13 chapitres × 2 pages $\approx$ 26 pages
\item \textbf{TOTAL:} $\approx$ 325 pages de contenu technique dense
\end{itemize}

\subsection*{Fichiers à Générer (hors optionnels)}

\begin{itemize}
\item 1 fichier TP-1
\item 1 fichier Ch-1 (Introduction)
\item Pour Ch0-Ch13 (14 chapitres): 5 fichiers chacun
\item \textbf{Total:} 1 + 1 + (14 × 5) = \textbf{72 fichiers}
\end{itemize}

\end{document}
